\section{Literature Review}
\label{sec:related}
% Draft for the authors to review and rewrite in their own words.
% Every claim below is checked against the literature matrix
% (paper_literature_and_course_checklist.xlsx, sheet "Literature matrix").

\subsection{Industrial anomaly detection and MVTec AD}
MVTec AD~\cite{bergmann2021mvtec} contains defect-free training images and
annotated test images for 15 categories, five of which are textures:
carpet, grid, leather, tile and wood. Most recent methods are evaluated on
it with image-level and pixel-level area under the receiver operating
characteristic curve (AUROC) and the per-region overlap (PRO). Because the
test pixels are overwhelmingly normal (about 97\% on MVTec
AD~\cite{li2024musc}), pixel AUROC is lenient, and studies increasingly
report pixel-level average precision (AP) and the area under the PRO curve
(AU-PRO) as well.

\subsection{Zero-shot single-image methods}
Aota et~al.~\cite{aota2023zeroshot} showed that texture anomalies can be
detected without any normal training images. Their method uses all patches
of the test image as the gallery and scores each patch by its mean distance
to its $K$ nearest neighbours ($K=400$ of 1{,}600 patches), using
WideResNet-50 features at $320\times320$ pixels. The method assumes that
normal regions form the majority of the image, and it is insensitive to the
orientation mismatch between training and test images that hurts many-shot
methods. FCA~\cite{ardelean2024highfidelity} replaces the nearest-neighbour
score with a comparison of local feature statistics against a median
reference, using a one-dimensional Wasserstein correspondence, and reports a
large reduction in localisation error on the MVTec AD textures; its
supplementary material notes failures when the anomaly covers the majority
of the image. Global-regularized neighborhood regression
(GRNR)~\cite{yao2025grnr} reconstructs each query from its neighbourhood with
a global normality prior, and NLPA-AD~\cite{zhang2025nlpa} propagates normal
labels across the texture while stopping propagation at likely anomalous
points. INP-Former~\cite{luo2025inp} learns to extract normal prototypes from
the test image itself, but it requires training. All of these methods take
their reference from the test image and report averages over all defects.

\subsection{Few-shot memory banks}
PatchCore~\cite{roth2022patchcore} introduced a memory bank of nominal patch
features from the full training set, scored by nearest-neighbour distance. AnomalyDINO~\cite{damm2025anomalydino}
showed that frozen DINOv2 patch features~\cite{oquab2024dinov2} make this
idea strong in the few-shot regime: with a memory bank built from a handful
of reference images and cosine distance to the nearest patch, it raised
one-shot image AUROC on MVTec AD from 93.1\% to 96.6\% without training.
SubspaceAD~\cite{lendering2026subspacead} models few-shot DINOv2 features
with a principal-component subspace and reports a PRO of about 97\% on the
textures with four references. Few-shot methods do not depend on the test
image being mostly normal, but they assume that the references are clean;
replacing three of ten references with defective images costs 1.8 AUROC points on average over four datasets (1.3 on MVTec AD)~\cite{villanueva2026cleaning}.

\subsection{Batched zero-shot methods}
MuSc~\cite{li2024musc} scores each patch against the patches of all other
unlabelled test images and aggregates local neighbourhoods at several scales,
small for small defects and large for large ones. Le-Gia and
Ahn~\cite{legia2025consistent} show that such methods fail when the same
anomaly recurs across test images, and they observe that these consistent
anomalies often have large masks. Batched methods need many test images at
once, which differs from the single-image setting studied here.

\subsection{Defect size and contamination}
Defect size has received little direct attention. Lehr
et~al.~\cite{lehr2024limitations} grow a defect pattern in steps and
measure the smallest reliably detected size for one method. The
contamination studied for few-shot detectors concerns defective images in
the reference set~\cite{villanueva2026cleaning}, not a defect inside the
test image that becomes part of its own reference.

\subsection{Summary and research gap}
Table~\ref{tab:lit} summarises where each family takes its reference. The
literature establishes that single-image methods assume a majority of
normal patches~\cite{aota2023zeroshot,ardelean2024highfidelity}, but it
reports results as averages over all defects. To our knowledge, no study has
compared a single-image and a few-shot reference across defect sizes on the
same features, or tested with an intervention whether contamination of the
image's own reference explains the difference. This is an empirical gap
(the size dependence has not been measured) and a methodological gap
(average scores cannot reveal it).

\begin{table}[t]
\centering
\caption{Where training-free texture anomaly detectors take their reference
for ``normal''.}
\label{tab:lit}
\footnotesize
\setlength{\tabcolsep}{3pt}
\begin{tabular}{>{\raggedright\arraybackslash}p{2.0cm}>{\raggedright\arraybackslash}p{1.5cm}>{\raggedright\arraybackslash}p{4.3cm}}
\toprule
Family & Examples & Reference and assumption \\
\midrule
Zero-shot, single image & \cite{aota2023zeroshot,ardelean2024highfidelity,yao2025grnr,zhang2025nlpa}
  & Rest of the test image; most of the image is normal \\
Memory bank (few-shot) & \cite{roth2022patchcore,damm2025anomalydino,lendering2026subspacead}
  & Patches of a few clean images; references are clean \\
Batched zero-shot & \cite{li2024musc,legia2025consistent}
  & Other unlabelled test images; anomalies do not recur \\
\bottomrule
\end{tabular}
\end{table}
